\section{Seam-count area law on the source family}
\label{sec:crossing-read-area-law}

The layered read law of the conservative source family supplies a finite
count that scales with area. Fix a population of constant density $n$,
sites per unit volume, in a bounded convex window
$\Omega\subset\mathbb R^3$, a read radius $a$, and the complete-neighbour
read law: an event at layer $j+1$ reads every event at layer $j$ within
distance $a$, the same site included. For a plane section
$\Sigma=\Omega\cap\{x_1=P\}$, a crossing read of one layer is an unordered
site pair $\{s,t\}$ with $s\neq t$, $\|t-s\|\leq a$, and $s_1<P<t_1$. Each
such pair is realized in every layer by one read in each direction across
$\Sigma$, from the lower layer to the upper layer; the ordered read count
is therefore twice the unordered count, and the unordered count per layer
is the primary quantity here. A same-site read never crosses a plane, and
a site with $s_1=P$ belongs to neither side.

\begin{proposition}[Crossing reads per unit area]
\label{prop:crossing-reads-per-area}
Let $S_k\subset\Omega$ be finite populations with a partition of $\Omega$
into cells of common volume $1/n_k$, each cell contained in the ball of
radius $H_k$ about its site, with read radii $a_k\to0$ and
$H_k/a_k\to0$. Let $\Sigma=\Omega\cap\{x_1=P\}$ have positive area
$|\Sigma|$. The number $X_k$ of unordered crossing reads of $\Sigma$ per
layer satisfies
\begin{equation}
 X_k=n_k^2\,|\Sigma|\int_0^{a_k}\frac{\pi(a_k-u)^2(2a_k+u)}{3}\,\dd u\,
     \bigl(1+o(1)\bigr)
    =\frac{\pi}{4}\,n_k^2a_k^4\,|\Sigma|\,\bigl(1+o(1)\bigr).
 \label{eq:crossing-read-density}
\end{equation}
The relative error is bounded by a constant multiple of $H_k/a_k+a_k$,
the constant depending on $\Omega$ and $\Sigma$.
\end{proposition}
\begin{proof}
Write $u(s)=P-s_1$ for the depth of a site below $\Sigma$ and
$C(s)=\{t\in\Omega:\ t_1>P,\ \|t-s\|\leq a_k\}$, so that
$X_k=\sum_{0<u(s)\leq a_k}|S_k\cap C(s)|$. The cell of a site $t$ lies in
$B(t,H_k)$; a cell that meets both $C(s)$ and its complement lies in the
$2H_k$-collar of $\partial C(s)$. Hence $|S_k\cap C(s)|$ differs from
$n_k\operatorname{vol}C(s)$ by at most $n_k$ times the volume of that
collar. The boundary $\partial C(s)$ consists of a spherical cap, a planar
disc and pieces of $\partial\Omega$, of total area $O(a_k^2)$, so the
error is $O(n_kH_ka_k^2)$ per site. The slab $0<u(s)\leq a_k$ contains
$O(n_ka_k|\Sigma|)$ sites, which gives a total error $O(n_k^2H_ka_k^3|\Sigma|)$,
of relative order $H_k/a_k$.

The function $f(s)=n_k\operatorname{vol}C(s)$ on the slab is bounded by
$\frac{2\pi}{3}n_ka_k^3$ and is Lipschitz with constant $O(n_ka_k^2)$: its
depth derivative is $-\pi n_k(a_k^2-u^2)$ and a lateral displacement changes
$C(s)$ through the faces of $\Omega$ by a set of measure at most $\pi a_k^2$
per unit displacement. The same cell argument, applied to the slab, bounds
$\bigl|\sum_{0<u(s)\leq a_k}f(s)-n_k\int_{0<u\leq a_k}f\,\dd^3s\bigr|$ by
$n_k$ times the Lipschitz constant times $H_k$ times the slab volume, plus
the contribution of the sites within $H_k$ of the two slab faces; both are
$O(n_k^2H_ka_k^3|\Sigma|)$, again of relative order $H_k/a_k$.

Finally, for $s$ at lateral distance more than $a_k$ from the relative
boundary of $\Sigma$ in the plane, $C(s)$ is the untruncated spherical cap
of $B(s,a_k)$ beyond height $u$, with volume
\[
 \int_u^{a_k}\pi(a_k^2-z^2)\,\dd z=\frac{\pi(a_k-u)^2(2a_k+u)}{3}.
\]
The truncated sites occupy a lateral collar of area $O(a_k)$ relative to
$|\Sigma|$, so the integral over the slab equals
$|\Sigma|\int_0^{a_k}\pi(a_k-u)^2(2a_k+u)/3\,\dd u$ up to a relative error
$O(a_k)$. The depth integral is $\pi a_k^4/4$.
\end{proof}

\begin{corollary}[Golden-family value]
\label{cor:golden-crossing-count}
For the golden source population $S_q$ of $q^3$ sites in
$\Omega=[0,L]^3$, $L=2/\sqrt{\varphi+2}$, with $n_q=q^3/L^3$,
$a_q=L/\sqrt q$ and cell assignment error $H_q=2\sqrt3L/q$, one has
$n_q^2a_q^4=q^4/L^2$, and the unordered crossing count per layer through
a full cross-section $\Sigma=\Omega\cap\{x_i=P\}$ is
\begin{equation}
 X_q=\frac{\pi q^4}{4}\bigl(1+O(q^{-1/2})\bigr).
 \label{eq:golden-crossing-count}
\end{equation}
For the mid-plane $P=L/2$ the continuum count with the caps truncated by
the cube is
\begin{equation}
 n_q^2\int_{s\in\Omega,\,s_1<L/2}
 \operatorname{vol}\{t\in\Omega:\ t_1>L/2,\ \|t-s\|\leq a_q\}\,\dd^3s
 =q^4\left(\frac{\pi}{4}-\frac{8}{15\sqrt q}+\frac1{12\,q}\right).
 \label{eq:golden-box-corrected}
\end{equation}
\end{corollary}
\begin{proof}
Here $H_q/a_q=2\sqrt3\,q^{-1/2}$ and $a_q/L=q^{-1/2}$, so
Proposition~\ref{prop:crossing-reads-per-area} gives
\eqref{eq:golden-crossing-count}. For \eqref{eq:golden-box-corrected}
substitute $d=t-s$. For $0<d_1\leq a_q\leq L/2$ and $|d_2|,|d_3|\leq L$,
the set of $s\in\Omega$ with $s+d\in\Omega$ and $s_1<L/2<s_1+d_1$ has
volume $d_1(L-|d_2|)(L-|d_3|)$. Over the half ball $d_1>0$,
$\|d\|\leq a_q$, the moments are $\int d_1=\pi a_q^4/4$,
$\int d_1|d_2|=4a_q^5/15$ and $\int d_1|d_2||d_3|=a_q^6/12$. Hence the
count equals $n_q^2a_q^4\bigl(\pi L^2/4-8La_q/15+a_q^2/12\bigr)$, which is
\eqref{eq:golden-box-corrected} after $n_q^2a_q^4=q^4/L^2$ and
$a_q=L/\sqrt q$.
\end{proof}

\begin{proposition}[Area law and the Planck-area dictionary]
\label{prop:planck-area-dictionary}
Declare that each crossing read carries one nat of cut entropy across
$\Sigma$ per layer. In the limit of
Proposition~\ref{prop:crossing-reads-per-area} the cut entropy is the area
law
\begin{equation}
 \Ent_{\mathrm{cut}}(\Sigma)=\frac{\pi}{4}\,n^2a^4\,|\Sigma| .
 \label{eq:seam-area-law}
\end{equation}
Identifying \eqref{eq:seam-area-law} with $|\Sigma|/(4\ell_P^2)$ fixes
\begin{equation}
 n^2a^4=\frac1{\pi\ell_P^2},\qquad
 \ell_P=\frac1{\sqrt\pi\,n\,a^2},
 \label{eq:planck-dictionary}
\end{equation}
and on the golden family $\ell_P=L/(\sqrt\pi\,q^2)$.
\end{proposition}
\begin{proof}
Equation \eqref{eq:seam-area-law} is \eqref{eq:crossing-read-density} with
one nat per read. Equating $(\pi/4)n^2a^4$ with $1/(4\ell_P^2)$ and
solving for $\ell_P^2$ gives \eqref{eq:planck-dictionary}; with
$n^2a^4=q^4/L^2$ this reads $\ell_P^2=L^2/(\pi q^4)$.
\end{proof}

\begin{remark}[What is not claimed]
\label{rem:crossing-not-claimed}
Equation \eqref{eq:planck-dictionary} is a calibration of the read scale by
the area coefficient, in the same way that the layer duration $a/c$ is a
calibration of the clock. No horizon is constructed: $\Sigma$ is a plane
section of the source window, with the same read law on both sides. The
nat-per-read reading is a declared dictionary, without a derivation from
the record channel of the finite ledger. No temperature, no Hawking flux
and no physical value of $\ell_P$ follow; $L$ is the side of the source
window, and the family fixes only the ratio $\ell_P/L=1/(\sqrt\pi\,q^2)$.
Dou and Sorkin~\cite{dousorkin2003} count causal links across a horizon in
a Poisson-sprinkled causal set and obtain an area law whose coefficient
depends on the sprinkling; here the population is deterministic, the count
is exact at every finite $q$, and the coefficient $\pi/4$ in
\eqref{eq:crossing-read-density} is the exact limit.
\end{remark}

The depth integral in \eqref{eq:crossing-read-density}, the value
$(\pi/4)n^2a^4$, the golden identity $n_q^2a_q^4=q^4/L^2$ with the
cross-section value $\pi q^4/4$, and the dictionary
\eqref{eq:planck-dictionary} with its golden value $\ell_P^2=L^2/(\pi q^4)$
are machine-checked in \texttt{Geometry/CrossingReadAreaLaw.lean}
(\texttt{capVolume\_integral}, \texttt{crossingDensity\_eq},
\texttt{golden\_crossing\_count}, \texttt{planck\_area\_of\_area\_law},
\texttt{golden\_planck\_length\_sq}) with the standard axioms only. The
limit statement of Proposition~\ref{prop:crossing-reads-per-area} is
analytic and is proved above.

\paragraph{Finite receipt.}
The exact crossing counts of the golden family at $q=5,8,13,21,34,55$ are
recorded with an independent verifier. The population and the neighbour
relation are rebuilt from the definitions, the relation agrees edge for
edge with the causal-poset evidence package at every level, every side
decision is exact in $\mathbb Q(\varphi)$, and no site lies on any of the
five planes $x_1=L/2$, $x_1=L/4$, $x_1=3L/4$, $x_2=L/2$, $x_3=L/2$.
Table~\ref{tab:crossing-receipt} lists the mid-plane counts. The counts
across $x_2=L/2$ and $x_3=L/2$ equal the $x_1=L/2$ count at every level,
as the product population is symmetric under coordinate permutations. At
$q=34$ the mid-plane count is $907\,080$ against the box-corrected value
$930\,602.0$, ratio $0.975$, and against the infinite-plane value
$1\,049\,555.8$, ratio $0.864$; at $q=55$ the count is $6\,496\,004$
against $6\,542\,684.6$, ratio $0.993$, and against $7\,186\,884.1$, ratio
$0.904$. The ratio to the box-corrected value is closer to one than the
ratio to the infinite-plane value at every level; it lies within $0.13$ of
one at every level and within $0.026$ of one for $q\geq21$, while the
ratio to the infinite-plane value rises from $0.658$ to $0.904$. The
residual at fixed $q$ is the finite population itself: the two sides of
the mid-plane hold $84\,700$ and $81\,675$ sites at $q=55$, and the mean
number of neighbours of a site, $1459.3$ at $q=55$, compares with the
cube mean $1461.1$ of $n_q\operatorname{vol}(B(s,a_q)\cap\Omega)$ and with
the untruncated ball value $1708.6$.

\begin{table}[ht]
\centering
\begin{tabular}{rrrrrr}
\toprule
$q$ & count & $\pi q^4/4$ & ratio & box-corrected & ratio \\
\midrule
5 & 323 & 490.9 & 0.658 & 352.2 & 0.917 \\
8 & 2\,164 & 3\,217.0 & 0.673 & 2\,487.3 & 0.870 \\
13 & 17\,754 & 22\,431.8 & 0.791 & 18\,390.1 & 0.965 \\
21 & 129\,830 & 152\,745.0 & 0.850 & 130\,882.5 & 0.992 \\
34 & 907\,080 & 1\,049\,555.8 & 0.864 & 930\,602.0 & 0.975 \\
55 & 6\,496\,004 & 7\,186\,884.1 & 0.904 & 6\,542\,684.6 & 0.993 \\
\bottomrule
\end{tabular}
\caption{Unordered crossing reads per layer across the mid-plane
$x_1=L/2$ of the golden family, the infinite-plane value $\pi q^4/4$, and
the box-corrected value \eqref{eq:golden-box-corrected}, with the ratios
of the count to each.}
\label{tab:crossing-receipt}
\end{table}
